\begin{table*}[t]
\centering
\small
\caption{Selected $a_0$ values (\AA{}). Experimental references follow Ref.~\citenum{Goldzak2022Solids}. BN and C reuse the AE curve for HD and HOC. These entries are not independent calculations.}
\label{tab:selected-a0_A}
\begin{tabular}{lrrrrr}
\toprule
Solid & Experiment & GX & AE & HD & HOC \\
\midrule
AlN & 4.3680 & 4.2811 & 4.3025 & 4.2747 & 4.2746 \\
AlP & 5.4480 & 5.3622 & 5.3781 & 5.3616 & 5.3621 \\
BN & 3.5930 & 3.5777 & 3.5675 & 3.5675 & 3.5675 \\
BP & 4.5250 & 4.4439 & 4.4696 & 4.4527 & 4.4527 \\
C & 3.5530 & 3.5345 & 3.5292 & 3.5292 & 3.5292 \\
LiCl & 5.0720 & 4.9081 & 4.9155 & 4.9178 & 4.9179 \\
MgO & 4.1890 & 4.1059 & 4.0935 & 4.0923 & 4.0919 \\
MgS & 5.1880 & 5.0689 & 5.0536 & 5.0700 & 5.0697 \\
Si & 5.4210 & 5.3386 & 5.3740 & 5.3383 & 5.3383 \\
SiC & 4.3470 & 4.2517 & 4.2920 & 4.2556 & 4.2556 \\
\bottomrule
\end{tabular}
\end{table*}

\begin{table*}[t]
\centering
\small
\caption{Selected $B_0$ values (GPa). Experimental references follow Ref.~\citenum{Goldzak2022Solids}. BN and C reuse the AE curve for HD and HOC. These entries are not independent calculations.}
\label{tab:selected-B0_GPa}
\begin{tabular}{lrrrrr}
\toprule
Solid & Experiment & GX & AE & HD & HOC \\
\midrule
AlN & 206.00 & 239.07 & 225.89 & 240.42 & 240.40 \\
AlP & 87.00 & 94.34 & 101.35 & 96.36 & 93.12 \\
BN & 388.50 & 423.39 & 420.87 & 420.87 & 420.87 \\
BP & 176.50 & 187.93 & 183.64 & 185.86 & 185.86 \\
C & 453.30 & 457.33 & 491.35 & 491.35 & 491.35 \\
LiCl & 37.30 & 47.98 & 46.37 & 47.13 & 47.16 \\
MgO & 173.00 & 188.28 & 198.71 & 213.13 & 207.37 \\
MgS & 81.00 & 87.90 & 90.27 & 84.17 & 86.22 \\
Si & 100.30 & 92.76 & 100.56 & 94.13 & 94.13 \\
SiC & 228.90 & 263.89 & 241.68 & 263.73 & 263.71 \\
\bottomrule
\end{tabular}
\end{table*}

\begin{table*}[t]
\centering
\small
\caption{EOS fit diagnostics, part 1. All full-window minima are bracketed. RMS residuals are in meV per atom.}
\label{tab:eos-fit-1}
\begin{tabular}{lrrrr}
\toprule
Method & Solid & $N_V$ & $B'_0$ & RMS \\
\midrule
AE & AlN & 9 & 3.707 & 0.0418 \\
AE & AlP & 9 & 4.179 & 0.0880 \\
AE & BN & 10 & 3.772 & 0.1250 \\
AE & BP & 10 & 3.701 & 0.0311 \\
AE & C & 10 & 3.756 & 0.0542 \\
AE & LiCl & 10 & 4.029 & 0.0205 \\
AE & MgO & 10 & 4.092 & 0.0190 \\
AE & MgS & 10 & 3.682 & 0.0117 \\
AE & Si & 10 & 4.544 & 0.3496 \\
AE & SiC & 10 & 3.854 & 0.0449 \\
GX & AlN & 9 & 2.802 & 0.6739 \\
GX & AlP & 9 & 3.959 & 0.1431 \\
GX & BN & 10 & 3.713 & 0.1306 \\
GX & BP & 10 & 4.081 & 0.0723 \\
GX & C & 10 & 3.611 & 0.1734 \\
GX & LiCl & 10 & 4.027 & 0.0435 \\
GX & MgO & 10 & 3.864 & 0.2921 \\
GX & MgS & 10 & 4.349 & 0.1597 \\
\bottomrule
\end{tabular}
\end{table*}

\begin{table*}[t]
\centering
\small
\caption{EOS fit diagnostics, part 2. All full-window minima are bracketed. RMS residuals are in meV per atom.}
\label{tab:eos-fit-2}
\begin{tabular}{lrrrr}
\toprule
Method & Solid & $N_V$ & $B'_0$ & RMS \\
\midrule
GX & Si & 10 & 5.096 & 0.5702 \\
GX & SiC & 10 & 4.107 & 0.1309 \\
HD & AlN & 9 & 3.339 & 0.3028 \\
HD & AlP & 9 & 4.778 & 0.3362 \\
HD & BP & 10 & 3.918 & 0.0532 \\
HD & LiCl & 10 & 3.905 & 0.0330 \\
HD & MgO & 10 & 4.208 & 0.3562 \\
HD & MgS & 10 & 3.167 & 0.2585 \\
HD & Si & 10 & 5.257 & 0.3634 \\
HD & SiC & 10 & 4.168 & 0.0561 \\
HOC & AlN & 9 & 3.338 & 0.3027 \\
HOC & AlP & 9 & 3.787 & 0.2878 \\
HOC & BP & 10 & 3.918 & 0.0532 \\
HOC & LiCl & 10 & 3.915 & 0.0328 \\
HOC & MgO & 10 & 4.181 & 0.0393 \\
HOC & MgS & 10 & 3.865 & 0.0993 \\
HOC & Si & 10 & 5.257 & 0.3634 \\
HOC & SiC & 10 & 4.167 & 0.0557 \\
\bottomrule
\end{tabular}
\end{table*}

\begin{table*}[t]
\centering
\small
\caption{Matched comparison with Ref.~\citenum{Goldzak2022Solids}: ten-solid set. Every row uses the identical intersection and the Goldzak static-lattice experimental reference. MAE($a_0$) is in \AA{} and MAE($B_0$) in GPa. A dash denotes unavailable data, not zero error. Native rows retain the numerical qualifications discussed in the text.}
\label{tab:literature-Goldzak2022}
\begin{tabular}{lrrr}
\toprule
Method & $N$ & MAE($a_0$) & MAE($B_0$) \\
\midrule
GX & 10 & 0.0831 & 16.61 \\
AE & 10 & 0.0728 & 16.89 \\
HD & 10 & 0.0844 & 21.77 \\
HOC & 10 & 0.0844 & 21.07 \\
HF & 10 & 0.0563 & 15.23 \\
MP2 & 10 & 0.0201 & 5.71 \\
SCS-MP2 & 10 & 0.0125 & 4.10 \\
SOS-MP2 & 10 & 0.0149 & 3.88 \\
\bottomrule
\end{tabular}
\end{table*}

\begin{table*}[t]
\centering
\small
\caption{Matched comparison with Ref.~\citenum{Zhang2018Solids}: nine solids, excluding AlN. Every row uses the identical intersection and the Goldzak static-lattice experimental reference. MAE($a_0$) is in \AA{} and MAE($B_0$) in GPa. A dash denotes unavailable data, not zero error. Native rows retain the numerical qualifications discussed in the text.}
\label{tab:literature-Zhang2018}
\begin{tabular}{lrrr}
\toprule
Method & $N$ & MAE($a_0$) & MAE($B_0$) \\
\midrule
GX & 9 & 0.0827 & 14.78 \\
AE & 9 & 0.0737 & 16.56 \\
HD & 9 & 0.0834 & 20.36 \\
HOC & 9 & 0.0835 & 19.59 \\
LDA & 9 & 0.0324 & 4.66 \\
PBE & 9 & 0.0451 & 12.96 \\
PBEsol & 9 & 0.0113 & 4.72 \\
M06-L & 9 & 0.0128 & 3.46 \\
SCAN & 9 & 0.0123 & 3.08 \\
HSE06 & 9 & 0.0170 & 5.86 \\
\bottomrule
\end{tabular}
\end{table*}

\begin{table*}[t]
\centering
\small
\caption{Matched comparison with Ref.~\citenum{Schimka2011HSEsol}: nine solids, excluding MgS. Every row uses the identical intersection and the Goldzak static-lattice experimental reference. MAE($a_0$) is in \AA{} and MAE($B_0$) in GPa. A dash denotes unavailable data, not zero error. Native rows retain the numerical qualifications discussed in the text.}
\label{tab:literature-Schimka2011}
\begin{tabular}{lrrr}
\toprule
Method & $N$ & MAE($a_0$) & MAE($B_0$) \\
\midrule
GX & 9 & 0.0791 & -- \\
AE & 9 & 0.0660 & -- \\
HD & 9 & 0.0807 & -- \\
HOC & 9 & 0.0807 & -- \\
PBE & 9 & 0.0438 & -- \\
PBEsol & 9 & 0.0136 & -- \\
HSE06 & 9 & 0.0132 & -- \\
HSEsol & 9 & 0.0117 & -- \\
\bottomrule
\end{tabular}
\end{table*}

\begin{table*}[t]
\centering
\small
\caption{Matched comparison with Ref.~\citenum{Mo2017Solids}: eight solids, excluding AlN and BN. Every row uses the identical intersection and the Goldzak static-lattice experimental reference. MAE($a_0$) is in \AA{} and MAE($B_0$) in GPa. A dash denotes unavailable data, not zero error. Native rows retain the numerical qualifications discussed in the text.}
\label{tab:literature-Mo2017}
\begin{tabular}{lrrr}
\toprule
Method & $N$ & MAE($a_0$) & MAE($B_0$) \\
\midrule
GX & 8 & 0.0911 & 12.27 \\
AE & 8 & 0.0797 & 14.58 \\
HD & 8 & 0.0907 & 18.86 \\
HOC & 8 & 0.0907 & 18.00 \\
Tao-Mo & 8 & 0.0194 & 4.26 \\
\bottomrule
\end{tabular}
\end{table*}

\begin{table*}[t]
\centering
\small
\caption{Matched comparison with Ref.~\citenum{Alizadeh2026PeriodicGFN2}: nine solids, excluding MgO. Every row uses the identical intersection and the Goldzak static-lattice experimental reference. MAE($a_0$) is in \AA{} and MAE($B_0$) in GPa. A dash denotes unavailable data, not zero error. Native rows retain the numerical qualifications discussed in the text.}
\label{tab:literature-PeriodicGFN2Manuscript}
\begin{tabular}{lrrr}
\toprule
Method & $N$ & MAE($a_0$) & MAE($B_0$) \\
\midrule
GX & 9 & 0.0831 & -- \\
AE & 9 & 0.0703 & -- \\
HD & 9 & 0.0831 & -- \\
HOC & 9 & 0.0830 & -- \\
GFN1-xTB & 9 & 0.1360 & -- \\
GFN2-xTB & 9 & 0.0659 & -- \\
\bottomrule
\end{tabular}
\end{table*}

\begin{table*}[t]
\centering
\small
\caption{Actual Gamma-centered k meshes by volume index. Indices 01--10 refer to the original evenly spaced volume ratios 0.90--1.10. Each selected volume uses the same mesh in all representations, but the mesh can change along an EOS.}
\label{tab:actual-k-meshes}
\begin{tabular}{lr}
\toprule
Solid & Mesh and volume indices \\
\midrule
AlN & $5^3$: 01,02,03,04,05,06,07,08,09 \\
AlP & $4^3$: 01,02,03,04,05,06,07,08,09 \\
BN & $6^3$: 02,03,04,05,06,07,08,09,10; $7^3$: 01 \\
BP & $5^3$: 01,02,03,04,05,06,07,08,09,10 \\
C & $6^3$: 04,05,06,07,08,09,10; $7^3$: 01,02,03 \\
LiCl & $5^3$: 01,02,03,04,05,06,07,08,09,10 \\
MgO & $5^3$: 06,07,08,09,10; $6^3$: 01,02,03,04,05 \\
MgS & $4^3$: 07,08,09,10; $5^3$: 01,02,03,04,05,06 \\
Si & $4^3$: 02,03,04,05,06,07,08,09,10; $5^3$: 01 \\
SiC & $5^3$: 01,02,03,04,05,06,07,08,09,10 \\
\bottomrule
\end{tabular}
\end{table*}
